\documentclass[11pt,a4paper]{article}
\usepackage[margin=25mm]{geometry}
\usepackage{amsmath,amssymb,amsthm,booktabs,hyperref}
\newtheorem{theorem}{Theorem}\newtheorem{proposition}[theorem]{Proposition}
\title{Even transitive markings on eleven points}
\author{Human Compression project}\date{}
\begin{document}\maketitle
We apply the complete companion paper \emph{Finite marked subgroups
by successive stabilizers}. Our target is Fernandes--Piedade--Reade,
\emph{String C-group representations of transitive groups: a case
study with degree 11}, arXiv:2302.11943v3, 1 November 2023, Theorem
8.8. Their proof uses permutation graphs and fracture graphs.
The proof here uses a complete finite enumeration with a written
coverage argument. It also recovers their principal permutation
marking list in Proposition 8.4 and its ten orbit-size entries.

\begin{theorem}[Fernandes--Piedade--Reade]\label{main}
Let $H\le A_{11}$ be transitive and have a string C-group
representation of rank four or five. Its rank is four, its
group is isomorphic to $\operatorname{PSL}_2(11)$, and its
abstract regular polytope is the self-dual $11$-cell.
\end{theorem}
\begin{proof}
Search all C-markings in the specified ambient $A_{11}$ through
rank five, modulo simultaneous conjugation by $S_{11}$. The
companion theorem proves coverage and uniqueness, including
markings generating proper transitive subgroups. The necessary
completion test discards only prefixes whose remaining pool
cannot generate a transitive group. Thus every marking in the
statement must appear in the final transitive list.

The complete execution of \texttt{enumerate\_even.g} visits
\[
  1,\ 2,\ 78,\ 2634,\ 147,\ 20
\]
prefix classes at lengths zero through five, discarding $2497$
branches by the necessary transitivity test. There are $130$
transitive markings of rank three, exactly two of rank four,
and none of rank five. Both rank-four groups have order $660$,
are explicitly isomorphic to $\operatorname{PSL}_2(11)$, and
have type $\{3,5,3\}$. The full recorded tuples, their group
computations and the complete search are required evidence,
not optional tests against the predicted count.

The two rank-four tuples are equivalent after reversing one
and conjugating points; as abstract ordered markings they are
also isomorphic. The exact generator homomorphism sending either
marking to its reverse is bijective. Their common abstract
marking is the $11$-cell: it agrees with source Proposition
8.4(5), whose literal permutations are printed below. The
already acquired projective classification identifies this
marking with the conventional $11$-cell and supplies its
geometric interpretation. Its full common paper and application,
not only the theorem statement, remain retained inputs. The
group--polytope correspondence gives the claimed uniqueness
and abstract self-duality.
\end{proof}

\paragraph{Literal source markings.}
For a graphical statement there is an additional identification
obligation. Number vertices in each row of the source diagram
from left to right, proceeding from the top row to the bottom.
The five source markings on physical page 23 become
\begin{align*}
T_1={}&((23)(45)(67)(89),\ (12)(34)(78)(9\,10),\
 &\qquad (28)(39)(67)(10\,11)),\\
T_2={}&((12)(34)(67)(89),\ (23)(45)(78)(9\,10),\
 &\qquad (16)(27)(34)(10\,11)),\\
T_3={}&((12)(34)(89)(10\,11),\ (23)(45)(78)(9\,10),\
 &\qquad (1\,10)(2\,11)(34)(56)),\\
T_4={}&((12)(67)(89)(10\,11),\ (23)(45)(78)(9\,10),\
 &\qquad (16)(27)(34)(10\,11)),\\
T_5={}&((12)(56)(89)(10\,11),\ (23)(45)(78)(10\,11),\
 &\qquad (34)(56)(8\,10)(9\,11),\\
 &\qquad (47)(58)(69)(10\,11)).
\end{align*}
Rows 1,2,4 use five top and six bottom vertices; row 3 uses
six top and five bottom vertices; row 5 uses rows of sizes
six, three and two. Parentheses here denote disjoint cycles,
not multiplication of adjacent decimal digits.

\begin{proposition}[The source permutation list]
Up to point conjugation and reversal, the transitive
$\operatorname{PSL}_2(11)$ markings of ranks three and four
are exactly $T_1,\ldots,T_5$. Up to abstract marked-group
isomorphism and reversal they give four polytopes, of types
$\{5,5\}$, $\{5,6\}$, $\{6,6\}$ and $\{3,5,3\}$.
\end{proposition}
\begin{proof}
The complete transitive rank-three census consists of $122$
markings generating $A_{11}$ and eight generating groups of
order $660$, explicitly identified with $\operatorname{PSL}_2(11)$.
The latter eight form four reversal pairs under point conjugation.
The rank-four list forms one such pair. Direct comparison with
the five displayed source tuples identifies all five classes,
with no additional class. A separate full-subset intersection
test checks each source tuple and every recorded transitive tuple.

For abstract equivalence, generator homomorphisms show that
$T_2$ and $T_3$ represent the same ordered marking, whereas
their point representations differ. The other three markings
are abstractly self-dual; $T_2,T_3$ have type $\{5,6\}$ and
are interchanged with the ordered dual type $\{6,5\}$ on
reversal. Distinct unordered types separate the four abstract
polytopes. All these homomorphism and conjugacy comparisons
are retained in the complete finite evidence.
\end{proof}

Deleting the first mark, and then doing the same after reversal,
gives the following orbit sizes, exactly source physical page 24:
\begin{center}
\begin{tabular}{ccc}\toprule
Source row & Suffix orbits & Reversed suffix orbits\\\midrule
1 & $1,5,5$ & $1,5,5$\\
2 & $2,3,6$ & $1,5,5$\\
3 & $2,3,6$ & $1,5,5$\\
4 & $2,3,6$ & $2,3,6$\\
5 & $1,10$ & $5,6$\\\bottomrule
\end{tabular}
\end{center}
The distinction between abstract and permutation self-duality
is essential: none of these five displayed permutation markings
is conjugate in $S_{11}$ to its reverse, although three underlying
abstract markings are self-dual.

\paragraph{Scope and attribution.}
Theorem~\ref{main} is the source's main result, not a new conjecture
answer. The rank-three and permutation classifications retain
the source's attribution to the earlier regular-map atlases;
the $11$-cell retains its attribution to Gr\"unbaum and Coxeter.
The complete finite search replaces the source's case analysis
for the stated principal contract. It does not assert recovery
of every independent auxiliary statement about arbitrary
permutation graphs, nor of the source appendix's $52$ claimed
independent generating tuples. These are separately disposed
in the source map, and are not premises of this classification.
No claim of literal theorem-set equivalence, numerical compression
or new Lean formalization is made.
\end{document}
